\documentclass[UTF8]{ctexart}
\usepackage{xeCJK}
\usepackage{geometry}
\geometry{left=2cm,right=2cm,top=2.5cm,bottom=2.5cm}
\setlength{\headheight}{12.7pt}

\usepackage{titlesec}
\titleformat{\section}
  {\normalfont\large\bfseries}
  {\thesection}
  {1em}
  {}
\titlespacing*{\section}
  {0pt}
  {3.5ex plus 1ex minus .2ex}
  {2.3ex plus .2ex}

\usepackage{amsmath}
\usepackage{amssymb}
\usepackage{amsthm}
\usepackage{enumitem}
\setlist[enumerate]{itemsep=0pt,topsep=0pt,parsep=0pt,leftmargin=0pt,itemindent=2.5em}
\setlist[itemize]{itemsep=0pt,topsep=0pt,parsep=0pt,leftmargin=0pt,itemindent=2.5em}
\usepackage{fancyhdr}
\pagestyle{fancy}
\fancyhf{}
\usepackage{graphicx}
\usepackage{hyperref}
\usepackage{indentfirst}
\usepackage{lmodern}


\begin{document}
\setlength{\abovedisplayskip}{1pt}
\setlength{\belowdisplayskip}{1pt}

\section{序数加法}

\hypertarget{sec:定义1.1}{}
\noindent\textbf{定义1.1 (序数加法)}

给定\href{../02 序数/01 序数#nameddest=sec:定义1.1}{序数} $\alpha$,
递归定义类映射$\mathbf{Add}_\alpha:\mathbf{On}\to\mathbf{On}$如下:
\vspace{3pt}
\[
    \forall\beta\in\mathbf{On}:\quad
    \mathbf{Add}_\alpha(\beta)\overset{\mathrm{def}}{=}\left\{\begin{aligned}
        &\alpha,\quad&&\beta=0,\\[-3pt]
        &\mathbf{Add}_\alpha(\max\beta)^+,\quad&&
        \beta\text{\,\,是后继序数},\\[-3pt]
        &\sup\mathbf{Add}_\alpha[\beta],\quad&&
        \beta\text{\,\,是极限序数}.
    \end{aligned}\right.
\]

接下来, 对任意的序数$\alpha,\beta$, 定义
\[
    \alpha+\beta\overset{\mathrm{def}}{=}\mathbf{Add}_\alpha(\beta),
\]
称为\textbf{序数加法}.

\vspace{10pt}

\hypertarget{sec:定理1.1}{}
\noindent\textbf{定理1.1 (序数加法)}

任给序数$\alpha,\beta$,
\begin{enumerate}
    \item $\alpha+0=\alpha$, $\alpha+1=\alpha^+$,
    \item $\alpha+\beta^+=(\alpha+\beta)^+$,
    \item 如果$\beta$是极限序数,
    那么$\displaystyle\alpha+\beta=\sup_{\gamma<\beta}(\alpha+\gamma)$.
\end{enumerate}

\vspace{15pt}

显然\href{../../03 自然数/02 自然数算术/01 自然数的加法#nameddest=sec:定义1.1}
{自然数的加法}是序数加法的特例. 易证任给序数$\alpha$有$0+\alpha=\alpha$.

\vspace{10pt}

\hypertarget{sec:定理1.2}{}
\noindent\textbf{定理1.2}

任给序数$\alpha,\beta,\beta'$, 如果$\beta<\beta'$,
那么$\alpha+\beta<\alpha+\beta'$.

\noindent{\kaishu 证明.}

给定$\alpha,\beta$, 对$\beta'>\beta$归纳证明$\alpha+\beta<\alpha+\beta'$.
任取$\beta'>\beta$, 假设
\[
    \forall\gamma:\quad\beta<\gamma<\beta'\to\alpha+\beta<\alpha+\gamma.
\]

\begin{enumerate}
    \item 如果$\beta'$是后继序数, 即有$\beta'=\gamma^+$,
    那么有$\beta\leqslant\gamma<\beta'$, 依假设有
    \[
        \alpha+\beta\leqslant\alpha+\gamma<\alpha+\gamma^+=\alpha+\beta'.
    \]
    \item 如果$\beta'$是极限序数, 则$\beta^+<\beta'$, 从而
    \begin{equation*}
        \alpha+\beta<\alpha+\beta^+\leqslant
        \sup_{\gamma<\beta'}(\alpha+\gamma)=\alpha+\beta'.
    \tag*{\qedsymbol}\end{equation*}
\end{enumerate}

\vspace{10pt}

\hypertarget{sec:定理1.3}{}
\noindent\textbf{定理1.3}

任给序数$\alpha,\alpha',\beta$, 如果$\alpha<\alpha'$,
那么$\alpha+\beta\leqslant\alpha'+\beta$.

\noindent{\kaishu 证明.}

给定$\alpha<\alpha'$, 对$\beta$归纳证明$\alpha+\beta\leqslant\alpha'+\beta$.
任取$\beta$, 假设对任意的$\gamma<\beta$有$\alpha+\gamma\leqslant\alpha'+\gamma$.

\begin{enumerate}
    \item 如果$\beta=0$, 显然.
    \item 如果$\beta$是后继序数, 即有$\beta=\gamma^+$, 那么$\gamma<\beta$, 依假设有
    \[
        \alpha+\beta=(\alpha+\gamma)^+\leqslant(\alpha'+\gamma)^+=\alpha'+\beta.
        \\[1pt]
    \]
    \item 如果$\beta$是极限序数, 那么依假设有
    \vspace{-1pt}
    \begin{equation*}
        \alpha+\beta=\sup_{\gamma<\beta}(\alpha+\gamma)\leqslant
        \sup_{\gamma<\beta}(\alpha'+\gamma)=\alpha'+\beta.
    \tag*{\qedsymbol}\end{equation*}
\end{enumerate}





\newpage

\hypertarget{sec:定理1.4}{}
\noindent\textbf{定理1.4 (序数减法)}

任给序数$\alpha,\beta$, 如果$\beta\leqslant\alpha$,
那么存在唯一的序数$\varepsilon$使得$\alpha=\beta+\varepsilon$, 记作$\alpha-\beta$.

\noindent{\kaishu 证明.}

\noindent{\kaishu 存在性.}

取满足$\beta+\varepsilon\geqslant\alpha$的最小的$\varepsilon$,
即对任意的$\gamma<\varepsilon$有$\beta+\gamma<\alpha$.

\begin{enumerate}
    \item 如果$\varepsilon=0$, 则$\beta+\varepsilon=\beta\leqslant\alpha$.
    \item 如果$\varepsilon$是后继序数, 即有$\varepsilon=\gamma^+$,
    那么$\gamma<\varepsilon$, 依假设有
    \[
        \beta+\varepsilon=(\beta+\gamma)^+\leqslant\alpha.
    \]
    \item 如果$\varepsilon$是极限序数, 那么依假设有
    \[
        \beta+\varepsilon=\sup_{\gamma<\varepsilon}(\beta+\gamma)\leqslant\alpha.
    \]
\end{enumerate}
综上有$\beta+\varepsilon\leqslant\alpha$, 即$\beta+\varepsilon=\alpha$.

\noindent{\kaishu 唯一性.} 因为加法保序, 显然. \hfill $\qedsymbol$

\vspace{10pt}

\hypertarget{sec:定理1.5}{}
\noindent\textbf{定理1.5}

任给序数$\alpha,\beta$, 如果$\beta$是后继(极限)序数,
那么$\alpha+\beta$也是后继(极限)序数.

\noindent{\kaishu 证明.}

如果$\beta$是后继序数, 即有$\beta=\gamma^+$,
那么$\alpha+\beta=(\alpha+\gamma)^+$是后继序数.

\vspace{3pt}
如果$\beta$是极限序数,
那么对任意的$\displaystyle x<\alpha+\beta=\bigcup_{\gamma<\beta}(\alpha+\gamma)$,
存在$\gamma<\beta$使得
\vspace{3pt}
\[
    x<\alpha+\gamma\quad\implies\quad x^+<\alpha+\gamma^+
    \quad\implies\quad x^+<\alpha+\beta,
\]
所以$\alpha+\beta$是归纳集, 即是极限序数. \hfill $\qedsymbol$

\vspace{10pt}

\hypertarget{sec:定理1.6}{}
\noindent\textbf{定理1.6 (序数加法的结合律)}

任给序数$\alpha,\beta,\gamma$, 有$(\alpha+\beta)+\gamma=\alpha+(\beta+\gamma)$.

\noindent{\kaishu 证明.}

给定$\alpha,\beta$, 对$\gamma$归纳证明. 任取$\gamma$, 假设对任意的$\delta<\gamma$
有$(\alpha+\beta)+\delta=\alpha+(\beta+\delta)$.

\begin{enumerate}
    \item 如果$\gamma=0$, 显然$(\alpha+\beta)+0=\alpha+\beta=\alpha+(\beta+0)$.
    
    \vspace{3pt}
    \item 如果$\gamma$是后继序数, 即有$\gamma=\delta^+$, 那么$\delta<\gamma$,
    依假设有
    \[
        (\alpha+\beta)+\gamma=((\alpha+\beta)+\delta)^+=
        (\alpha+(\beta+\delta))^+=\alpha+(\beta+\delta)^+=\alpha+(\beta+\gamma).
    \]

    \vspace{3pt}
    \item 如果$\gamma$是极限序数, 此时$\beta+\gamma$也是极限序数.
    
    \vspace{3pt}\hspace{2.17em}
    一方面, 对任意的$\delta<\gamma$, 有
    \vspace{3pt}
    \[
        \alpha+(\beta+\delta)\leqslant
        \sup_{\theta<\beta+\gamma}(\alpha+\theta)\quad\implies\quad
        \sup_{\delta<\gamma}(\alpha+(\beta+\delta))\leqslant
        \sup_{\theta<\beta+\gamma}(\alpha+\theta).
    \]

    \vspace{3pt}\hspace{2.17em}
    另一方面, 对任意的$\displaystyle\theta<\beta+\gamma=
    \bigcup_{\delta<\gamma}(\beta+\delta)$, 存在$\delta<\gamma$使得
    $\theta<\beta+\delta$, 所以
    \vspace{3pt}
    \[
        \alpha+\theta<\alpha+(\beta+\delta)\leqslant\,
        \sup_{\delta<\gamma}(\alpha+(\beta+\delta))\quad\implies\quad
        \sup_{\theta<\beta+\gamma}(\alpha+\theta)\leqslant
        \sup_{\delta<\gamma}(\alpha+(\beta+\delta)).
    \]

    \hspace{2.17em}
    综上即得
    \begin{equation*}
        (\alpha+\beta)+\gamma=\sup_{\delta<\gamma}((\alpha+\beta)+\delta)=
        \sup_{\delta<\gamma}(\alpha+(\beta+\delta))=
        \sup_{\theta<\beta+\gamma}(\alpha+\theta)=\alpha+(\beta+\gamma).
    \tag*{\qedsymbol}\end{equation*}
\end{enumerate}

\end{document}